\documentclass[10pt,twocolumn]{article}
\usepackage[margin=0.72in]{geometry}
\usepackage{booktabs}
\usepackage{graphicx}
\usepackage{microtype}
\usepackage{enumitem}
\usepackage{xcolor}
\usepackage{hyperref}
\usepackage{amsmath}
\usepackage[numbers,sort&compress]{natbib}

\definecolor{bayanat}{HTML}{5B4DF2}
\hypersetup{colorlinks=true,linkcolor=bayanat,citecolor=bayanat,urlcolor=bayanat}
\setlist{nosep,leftmargin=*}
\setlength{\columnsep}{0.25in}
\setlength{\parindent}{1em}
\setlength{\parskip}{0pt}

\title{\vspace{-1.5em}\textbf{Arabic Annotation Is Not Binary:}\\
\large LabelBench, a Reproducible Audit of Automated Data Labeling}
\author{Bayanat Labs Research\\
\small \href{https://www.bayanatlabs.com}{bayanatlabs.com}}
\date{Technical Report v0.1, revised July 2026}

\begin{document}
\maketitle

\begin{abstract}
Language models are increasingly proposed as substitutes for human data
annotators, but aggregate accuracy on simple tasks can hide failure on the
taxonomies, dialects, and risk thresholds that make production Arabic labeling
difficult. We introduce \textsc{LabelBench}, a checksummed, resumable audit
harness for Arabic annotation. The first release covers binary toxicity,
seven-way multi-label toxicity, 60-way intent classification, and five-region
dialect identification. It reports exact-match accuracy, macro-F1, paired
bootstrap intervals, calibration diagnostics, source-leakage audits, and a
conservative selective-automation diagnostic: the fraction of items that can
be accepted while a Wilson lower bound remains above a target accuracy.

The headline finding is an annotation complexity cliff. On identical held-out
samples, Qwen 2.5 3B reaches 78.0\% accuracy on binary toxicity but only
25.3\% exact match on the full toxicity taxonomy. Arabic-specialized Fanar 1
9B improves those scores to 88.0\% and 31.3\%, respectively, while a small
character n-gram supervised diagnostic reaches 97.0\% and 49.0\%. On a
balanced native dialect set, Qwen 2.5 3B reaches 22.4\% accuracy, near the
20\% chance floor, and predicts Egyptian for 70 to 80\% of every non-Egyptian
region. Fanar improves dialect accuracy to 47.8\%, but remains 22.0 points
behind a 500-example character diagnostic. Because the v0.1 generative runs
emit categorical JSON rather than calibrated confidence, their zero selective
coverage is a harness limitation, not evidence that selective automation is
impossible. The contribution is therefore not a state-of-the-art model ranking.
It is a reproducible Arabic audit protocol showing how taxonomy, dialect,
confidence, and human review must be measured before automated labels are
trusted.
\end{abstract}

\section{Introduction}
The commercial promise of model-assisted annotation is straightforward: a
model proposes labels, humans review only uncertain cases, and cost falls
without sacrificing quality. For Arabic, that promise is tempting and
dangerous. Arabic production queues are not one language in one register.
They mix dialects, countries, scripts, code-switching, profanity conventions,
domain-specific policies, and label taxonomies whose operational meaning
depends on fine distinctions. A model that looks competent on a binary proxy
task can still fail at the decision actually needed by a moderation queue, a
retrieval evaluator, a speech system, or a post-training data engine.

This paper asks a narrower and more actionable question than ``Can models label
Arabic?'' We ask:
\begin{enumerate}
  \item When the label space becomes more realistic, how quickly does
  zero-shot automated labeling degrade?
  \item Does an Arabic-specialized local model reduce the failure modes of a
  generic local instruction model?
  \item Do the evaluated outputs include enough calibrated confidence signal to
  justify auto-accepting any portion of the queue?
  \item Can a laptop-scale benchmark expose these issues reproducibly without
  redistributing restricted text?
\end{enumerate}

The answer is sobering but useful. Specialization helps: Fanar 1 9B beats Qwen
2.5 3B by 10.0 points on binary toxicity, 6.0 points on multi-label exact
match, and 25.4 points on dialect identification. However, specialization does
not, under this zero-shot JSON-only protocol, provide a calibrated auto-accept
policy at 95\% accuracy. A small supervised diagnostic still outperforms the
evaluated zero-shot systems on all paired tasks. The narrow lesson is that
Arabic labeling automation requires task-specific measurement, calibrated
confidence, and native review rather than generic ``AI annotator'' claims.

\section{Contributions}
We make five contributions.
\begin{enumerate}
  \item \textbf{A reproducible Arabic annotation audit.} LabelBench pins
  source revisions, raw-file checksums, normalized manifest hashes, model
  digests, prompts, runtime metadata, and append-only prediction logs.
  \item \textbf{A production-oriented selective-automation diagnostic.} SAC is
  a simple Wilson-bound coverage calculation connected to selective
  classification, calibration, and conformal risk-control work
  \citep{geifman2017selective,guo2017calibration,angelopoulos2024conformal}.
  \item \textbf{A quantified source-confound audit.} We show that combined
  Arabic dialect corpora can let models learn corpus provenance rather than
  dialect; a trivial source-majority shortcut would reach 80.1\% on combined
  IADD.
  \item \textbf{Paired evidence for an annotation complexity cliff.} The same
  held-out items are used for Qwen, Fanar, and supervised diagnostics, enabling
  paired bootstrap intervals and exact McNemar tests.
  \item \textbf{A human adjudication path.} The release includes a blinded,
  right-to-left, offline review bundle so future versions can separate model
  errors from gold ambiguity and taxonomy defects.
\end{enumerate}

\section{Scope of the claim}
The paper is intentionally framed as an audit report rather than a leaderboard.
This distinction matters. A leaderboard asks which system is best under a
fixed benchmark. An annotation audit asks whether a system is safe to insert
into a labeling workflow with a stated quality target. The latter requires
item-level reproducibility, confidence calibration, gold-label scrutiny, and
an escalation rule for uncertain examples. A model can improve average
accuracy and still be unusable for automatic acceptance if its confidence
ranking does not separate correct from incorrect labels. Conversely, a model
with imperfect average accuracy can be operationally useful if a small,
well-calibrated slice can be accepted and the rest escalated.

LabelBench v0.1 therefore supports three limited claims. First, the selected
Arabic manifests expose a large gap between binary proxies and production-like
taxonomies. Second, the evaluated laptop-scale zero-shot systems are not
validated automation policies under the current prompting and confidence
protocol. Third, the harness is sufficient to run stronger tests: frontier
models, few-shot/native prompts, human adjudication, and risk-controlled
selection can be added without changing the provenance or logging contract.
It does not support a general claim that all generative models fail at Arabic
annotation, nor a claim that character n-gram diagnostics are state of the art.

\section{Related work}
\paragraph{Arabic resources and shared tasks.}
AraTox provides a multi-dialect, multi-label Arabic toxicity benchmark with
expert labels across six Arabic varieties \citep{aratox2026}. MASSIVE provides
localized virtual-assistant utterances across many languages, including
ar-SA, and is useful as a high-structure intent classification contrast
\citep{massive2022}. IADD integrates multiple dialect identification datasets
\citep{iadd2022}; DART is a crowd-labeled Arabic dialect tweet corpus covering
the five broad regions used here \citep{dart2018}. LabelBench is not a
replacement for community benchmarks: NADI is the canonical Arabic dialect-ID
shared-task series and now includes multi-label country-level dialect
identification and dialectness estimation \citep{abdulmageed2024nadi}. OSACT5
is the relevant shared-task lineage for Arabic offensive language and
fine-grained hate-speech detection \citep{mubarak2022osact5}. Future
LabelBench releases should add NADI/OSACT-compatible tasks so results are
directly comparable to the field.

\paragraph{Dialect ambiguity and dataset shortcuts.}
Arabic dialect identification is not a clean single-label problem. Keleg and
Magdy show that many apparent system errors reflect incomplete dialect labels
and argue for multi-label framing \citep{keleg2023scrutiny}. This motivates
our human adjudication protocol and limits how strongly any DART-only
five-region result should be interpreted.

\paragraph{Annotation quality.}
The data-labeling literature has long shown that cheap annotation can be
useful when disagreement, worker quality, and aggregation are measured rather
than assumed \citep{snow2008cheap,dawid1979maximum}. LabelBench adopts the
same posture toward automated annotators: a model prediction is not a label
until its quality is measured under the relevant taxonomy and review policy.
Recent LLM-annotation work shows that frontier models can sometimes match or
exceed crowd annotators on selected text tasks \citep{gilardi2023chatgpt}, that
zero-shot LLMs can assist computational social-science annotation while still
lagging fine-tuned classifiers on taxonomic labeling \citep{ziems2024large},
and that off-the-shelf multilingual LLM annotation varies substantially across
tasks and low-resource languages \citep{bhat2023large}. The Alternative
Annotator Test provides a statistical procedure for deciding when an LLM can
replace human annotations on a task-specific sample
\citep{calderon2025alttest}. LabelBench should be read as an audit harness
that can feed such procedures, not as a substitute for them.

\paragraph{Selective prediction and calibration.}
Selective classification evaluates the tradeoff between coverage and risk:
a system may abstain when it is likely to be wrong \citep{geifman2017selective}.
Calibration asks whether stated confidence reflects empirical correctness
\citep{guo2017calibration}. In annotation operations, these are not academic
extras. They determine whether a queue can be partially automated without
quietly corrupting training data. SAC is a deliberately simple launch
diagnostic; conformal risk control and selective conformal risk control provide
stronger finite-sample frameworks for future releases
\citep{angelopoulos2024conformal,xu2025selective}.

\section{Benchmark design}
\subsection{Tasks and data}
Table~\ref{tab:tasks} summarizes LabelBench v0.1. AraTox contains
approximately 36,700 expert-annotated comments spanning six Arabic varieties
and seven toxicity categories \citep{aratox2026}. We derive a binary toxicity
view and retain the original multi-label view. MASSIVE ar-SA is included as a
translated, high-structure intent benchmark rather than treated as native
dialect data. For dialect identification we isolate only the DART source inside
IADD, for the reasons in Section~\ref{sec:source-confound}.

\begin{table*}[t]
\centering
\small
\caption{LabelBench v0.1 tasks. Train counts apply only to supervised
diagnostics; zero-shot models receive no demonstrations.}
\label{tab:tasks}
\begin{tabular}{llllrr}
\toprule
Task & Source & Native / localized & Labels & Train & Test \\
\midrule
Binary toxicity & AraTox & native, multi-dialect & 2 & 33,038 & 3,670 \\
Multi-label toxicity & AraTox & native, multi-dialect & 7 & 33,038 & 3,670 \\
Intent classification & MASSIVE ar-SA & localized & 60 & 11,514 & 2,974 \\
Dialect region & IADD/DART only & native, crowd-labeled & 5 & 500 & 500 \\
\bottomrule
\end{tabular}
\end{table*}

\subsection{Why source confounding matters}
\label{sec:source-confound}
Dialect identification datasets often combine corpora collected for different
projects. If one corpus contributes mostly Levantine text and another mostly
Maghrebi text, a classifier can learn collection artifacts, topic patterns, or
annotation provenance instead of dialect. IADD is valuable because it
integrates several resources, but that integration also creates a serious
evaluation trap.

Table~\ref{tab:source} quantifies the issue. A trivial rule that predicts each
IADD source's majority region is 80.1\% accurate over the five target regions.
This number is higher than every dialect result in our DART-only evaluation,
including the supervised diagnostic. Reporting such a score as dialect
competence would be misleading. We therefore evaluate only DART, which contains
all five target regions under one source. After normalization, DART yields
1,761 clean unique rows: 420 Gulf, 419 Egyptian, 391 Levantine, 328 Maghrebi,
and 203 Iraqi. We draw a balanced 100-per-region evaluation set and a disjoint
100-per-region diagnostic training set.

\begin{table}[t]
\centering
\small
\caption{IADD source-region confounding. Counts are examples in the five
regions evaluated here. The source-majority shortcut would score 80.1\%.}
\label{tab:source}
\begin{tabular}{lrrrrrr}
\toprule
Source & EGY & GLF & IRQ & LEV & MGH & Total \\
\midrule
AOC & 4414 & 6259 & 9 & 6503 & 28 & 17213 \\
DART & 423 & 423 & 207 & 397 & 331 & 1781 \\
PADIC & 0 & 0 & 0 & 14426 & 21639 & 36065 \\
SHAMI & 0 & 0 & 0 & 66247 & 0 & 66247 \\
TSAC & 0 & 0 & 0 & 0 & 11998 & 11998 \\
\bottomrule
\end{tabular}
\end{table}

\subsection{Integrity controls}
Every network source is pinned to a content revision and SHA-256 digest. Every
normalized train, validation, test, and sampled manifest has a second digest.
Prediction rows are appended atomically and keyed by stable item ID; resumed
runs validate cached gold labels and reject unknown or duplicate IDs. Reports
record exact local model digests, prompt and completion counts where available,
runtime, Python version, code revision, and prediction-file digest.

We audit normalized train/test overlap before scoring. MASSIVE ar-SA contains
8.41\% normalized overlap and 63 conflicting overlaps; the character baseline
falls from 74.24\% to 73.27\% after overlapped test items are removed. AraTox
binary overlap is 0.49\% with no binary-label conflicts and no material score
change. We report these findings rather than silently deduplicating a canonical
test set.

\subsection{Models}
The generative baselines use locally quantized instruction models through
Ollama with temperature zero, a fixed seed, and strict JSON outputs. Qwen 2.5
3B is the generic laptop-scale principal baseline. Fanar 1 9B is the
Arabic-specialized local model. Qwen 2 7B and Llama 3 8B are included on the
dialect task to check whether the Qwen 2.5 dialect collapse is merely a small
model artifact. Structured-output validity is scored rather than repaired.
These are laptop-scale baselines, not the systems a well-funded production
team would necessarily deploy. Frontier and larger Arabic-specialist models
remain necessary before making broad model-capability claims.

Supervised diagnostics are intentionally simple: a character 2 to 5-gram
multinomial naive Bayes classifier for single-label tasks and a one-vs-rest
variant for multi-label toxicity. Multi-label thresholds are selected on a
15\% partition of the training split, then the models are refit on all
training data. These diagnostics are not state-of-the-art claims. They answer
a narrower question: does the held-out manifest contain learnable signal under
a small-compute regime?

\subsection{Prompting and confidence protocol}
The local generative runs use a single zero-shot instruction template per task.
The prompt names the closed label set, asks for strict JSON, and does not
include demonstrations, rationales, chain-of-thought, or Arabic-native prompt
variants. Decoding is deterministic. This design makes the first release
cheap, reproducible, and easy to inspect, but it also deliberately under-tests
what a production annotation team would try. Few-shot examples, Arabic prompts,
self-consistency, and model-specific output calibration are part of the next
experiment matrix.

The original v0.1 runs store categorical labels with schema-level confidence
set to 1.0. That value records parser validity, not model belief. The current
code adds a verbalized-confidence mode for future local runs and an OpenAI
Responses adapter that can require a confidence field in the JSON schema. Those
extensions are not retroactively applied to the v0.1 numbers. They are included
so that the next release can separate model accuracy from the routing question:
which items, if any, can be accepted automatically?

\section{Metrics}
For single-label tasks we report accuracy, macro-F1, a 1,000-replicate
bootstrap 95\% interval for accuracy, Brier score, expected calibration error,
and a confusion matrix. For multi-label toxicity the primary metric is subset
accuracy: a prediction is correct only when its entire label set matches.
Micro-F1, macro-F1, Hamming loss, label cardinality, and per-label F1 provide
secondary views.

\paragraph{Safe Automation Coverage.}
For predictions sorted by confidence, consider every threshold without
splitting ties. At each threshold, let $k$ be correct among $n$ accepted
predictions and let $L(k,n)$ be the 95\% Wilson lower confidence bound
\citep{wilson1927}. For quality target $\tau$,
\[
  \mathrm{SAC}_{\tau} =
  \max_t \frac{n_t}{N}\quad \mathrm{s.t.}\quad L(k_t,n_t) \geq \tau.
\]
We report SAC95 and SAC98 as an operational diagnostic, not as a new
statistical guarantee. Empirical accuracy alone is insufficient when the
accepted set is small, but selecting a threshold on the same evaluation set is
optimistic. The code now supports calibration-split SAC, where the threshold is
chosen on a calibration partition and evaluated separately. A future benchmark
release should replace or accompany SAC with conformal risk-control methods.
Models that emit only schema-level confidence of 1.0 cannot manufacture a
useful ranking and therefore must satisfy the bound on the whole tied set.

\paragraph{Paired tests.}
When systems share the same evaluation items, we report paired bootstrap
intervals for accuracy or exact-match differences and exact two-sided McNemar
tests on the correctness contingency table. The paired design is important:
it distinguishes a real model difference from an easier or harder sample.
For the eight central v0.1 paired comparisons, the release JSON also reports
Holm-adjusted McNemar $p$-values. Effects that do not survive this correction
are treated as exploratory, even when their raw paired tests are below 0.05.

\section{Results}
\begin{figure*}[t]
\centering
\includegraphics[width=0.88\textwidth]{figures/annotation-gap.pdf}
\caption{The annotation complexity cliff on identical held-out manifests.
Binary and multi-label toxicity use the same seeded AraTox samples for all
three systems. The dialect diagnostic is trained on 500 DART items disjoint
from the balanced 500-item evaluation manifest.}
\label{fig:gap}
\end{figure*}

Figure~\ref{fig:gap} gives the central result. Qwen's 78.0\% binary toxicity
accuracy drops to 25.3\% when exact agreement with the seven-label taxonomy is
required. Fanar improves the corresponding scores to 88.0\% and 31.3\%, but
does not close the gap to small supervised diagnostics at 97.0\% and 49.0\%.
For five-region dialect identification, Qwen scores 22.4\% accuracy and 14.9\%
macro-F1 against a 20\% chance floor. Fanar reaches 47.8\%, while a
500-example character diagnostic reaches 69.8\% accuracy and 68.8\%
macro-F1.

\begin{table*}[t]
\centering
\small
\caption{Held-out results. Exact is accuracy for single-label tasks and
subset accuracy for multi-label toxicity. SAC95 is the maximum coverage whose
Wilson lower bound remains at least 95\%.}
\label{tab:results}
\begin{tabular}{llrrrrr}
\toprule
Task & Model & Items & Exact & Macro-F1 & Micro-F1 & SAC95 \\
\midrule
Binary toxicity & Qwen 2.5 3B & 100 & 78.0 & 69.8 & -- & 0.0 \\
Binary toxicity & Fanar 1 9B & 100 & 88.0 & 83.1 & -- & 0.0 \\
Binary toxicity & Char NB & 100 & \textbf{97.0} & \textbf{95.0} & -- & 89.0 \\
Multi-label toxicity & Qwen 2.5 3B & 300 & 25.3 & 39.3 & 43.0 & 0.0 \\
Multi-label toxicity & Fanar 1 9B & 300 & 31.3 & 50.5 & 55.0 & 0.0 \\
Multi-label toxicity & Char OVR & 300 & \textbf{49.0} & \textbf{73.2} & \textbf{73.9} & 0.0 \\
Dialect region & Qwen 2.5 3B & 500 & 22.4 & 14.9 & -- & 0.0 \\
Dialect region & Qwen 2 7B & 500 & 22.6 & 16.2 & -- & 0.0 \\
Dialect region & Llama 3 8B & 500 & 24.6 & 18.7 & -- & 0.0 \\
Dialect region & Fanar 1 9B & 500 & 47.8 & 47.3 & -- & 0.0 \\
Dialect region & Char NB & 500 & \textbf{69.8} & \textbf{68.8} & -- & 0.0 \\
Intent classification & Char NB & 2974 & 74.2 & 65.1 & -- & 39.5 \\
\bottomrule
\end{tabular}
\end{table*}

\subsection{Binary toxicity is not enough}
Binary toxicity is the easiest and most forgiving task in the release. On the
100-item paired slice, Qwen reaches 78.0\% accuracy. Fanar reaches 88.0\%,
with only one false positive but eleven toxic items missed as non-toxic. The
character diagnostic reaches 97.0\%. On the full AraTox binary test set, the
character diagnostic reaches 96.27\% accuracy and 95.47\% macro-F1, while a
majority baseline is 71.88\% accurate but only 41.82\% macro-F1. This confirms
that the binary label is learnable and that raw accuracy can be inflated by
class imbalance.

\subsection{Multi-label toxicity exposes category errors}
The multi-label task is operationally closer to a real policy queue. A system
must choose among appearance attacks, cussing, hatred, racial attacks, sexual
content, violence, and the mutually exclusive ``not'' label. On the 300-item
paired sample, Qwen's exact match is 25.3\% and macro-F1 is 39.3\%. Fanar
improves to 31.3\% exact match and 50.5\% macro-F1. The tuned character OVR
diagnostic reaches 49.0\% exact match and 73.2\% macro-F1 on the same sample.
On the full AraTox test set, the tuned diagnostic reaches 53.73\% exact match,
73.35\% macro-F1, 74.70\% micro-F1, and 11.64\% Hamming loss.

\begin{table}[t]
\centering
\small
\caption{Per-label F1 on the 300-item multi-label toxicity sample. The weak
categories differ by model, which matters for policy routing.}
\label{tab:perlabel}
\begin{tabular}{lrrr}
\toprule
Label & Qwen & Fanar & Char OVR \\
\midrule
Appearance & 34.7 & 20.0 & 72.3 \\
Cussing & 26.2 & 62.9 & 73.6 \\
Hatred & 50.0 & 47.3 & 70.5 \\
Racial & 36.4 & 39.0 & 46.7 \\
Sexual & 49.5 & 71.8 & 86.9 \\
Violence & 10.3 & 42.3 & 70.6 \\
Not & 68.1 & 69.9 & 92.1 \\
\bottomrule
\end{tabular}
\end{table}

Table~\ref{tab:perlabel} shows that the failure is not merely low exact match.
Qwen is particularly weak on violence, with 10.3\% F1. Fanar raises violence
F1 to 42.3\% and sexual-content F1 to 71.8\%, but drops appearance F1 to
20.0\%. The average number of predicted labels is close to the gold average
for Qwen (1.40 predicted versus 1.38 gold), so the problem is not only
under-prediction. The model often chooses the wrong categories.

\subsection{Dialect collapse is systematic}
Qwen 2.5 predicts Egyptian for 85\% of Egyptian items, but also for 75\% of
Gulf, 70\% of Iraqi, 80\% of Levantine, and 70\% of Maghrebi items. Recall is
4\% for Gulf, 6\% for Iraqi, 13\% for Levantine, and 4\% for Maghrebi. Qwen 2
7B changes the dominant prediction from Egyptian toward Levantine but reaches
nearly identical accuracy. Llama 3 reaches 24.6\%. This indicates that the
failure is not just one small model being noisy.

Fanar 1 9B reaches 47.8\%, with 60\% Egyptian, 75\% Gulf, 13\% Iraqi, 50\%
Levantine, and 41\% Maghrebi recall. Regional specialization substantially
reduces collapse, although Iraqi remains a severe weakness and the supervised
diagnostic remains 22 points higher. The character diagnostic's recalls are
81\% Egyptian, 70\% Gulf, 88\% Iraqi, 70\% Levantine, and 40\% Maghrebi.
Maghrebi remains difficult even for the diagnostic, consistent with DART's
reported lower agreement for some dialect labels \citep{dart2018}.

\subsection{Paired significance tests}
All central comparisons use identical items. Exact McNemar $p$-values below
are raw unless otherwise stated; machine-readable results include Holm
family-wise corrections across the eight paired comparisons. Fanar exceeds
Qwen by 10.0 points on binary toxicity (paired-bootstrap 95\% CI: 2.0 to 18.0;
raw $p=0.031$, Holm $p=0.090$), while the character diagnostic exceeds Fanar
by 9.0 points (2.0 to 17.0; raw $p=0.035$, Holm $p=0.090$) and Qwen by 19.0
points (11.0 to 28.0; Holm $p=8.4\times10^{-5}$).

On multi-label toxicity, the character diagnostic exceeds Qwen by 23.7 points
(17.7 to 29.7; Holm $p<10^{-12}$). Fanar exceeds Qwen by 6.0 points (1.0 to
11.3; raw $p=0.030$, Holm $p=0.090$), while the character diagnostic remains
17.7 points higher than Fanar (11.7 to 23.7; Holm $p<5\times10^{-7}$). On
dialect identification, Fanar exceeds Qwen by 25.4 points (20.2 to 30.6; Holm
$p<2\times10^{-18}$), while the character diagnostic remains 22.0 points above
Fanar (16.2 to 27.6; Holm $p<2\times10^{-12}$).

\subsection{Confidence was not elicited, so SAC is diagnostic only}
The zero-shot structured classifiers return categorical schema outputs rather
than calibrated class probabilities, creating a single confidence tie. None
reaches non-zero SAC95 under this tied-confidence protocol. This should not be
read as evidence that generative models cannot support selective automation.
It shows that a JSON label alone is insufficient for automation routing. The
harness now supports a verbalized-confidence mode for future local-model runs;
frontier models should additionally be tested with token log probabilities,
self-consistency agreement, or conformal risk-control calibration where
available. The tuned multi-label character system also reaches zero SAC95
despite stronger F1, showing that label quality and safe triage are separate
problems. The only substantial SAC95 in v0.1 is the binary character
diagnostic: 89.0\% on the seeded 100-item sample and 100\% on the full AraTox
binary test. Intent classification reaches 39.5\% SAC95, a reminder that
selective automation is possible when confidence ranking is useful.

\section{Robustness checks}
The release is designed to reduce three common benchmark failure modes. First,
the DART-only dialect protocol avoids the 80.1\% source-majority shortcut in
combined IADD. Second, all central model comparisons use paired items, so
differences are tested on the same examples rather than across different
samples. Third, the smaller paired LLM slices are interpreted alongside
full-test supervised diagnostics, which confirm that the binary and
multi-label AraTox tasks are learnable under a low-compute text baseline. These
checks support a narrow claim: under one zero-shot local-model protocol,
including an Arabic-specialized model, these systems do not yet provide a
validated Arabic label-automation policy on these taxonomies without measured
human review.

\section{Human adjudication protocol}
Automated scores cannot detect every bad gold label, ambiguous dialectal form,
or taxonomy defect. LabelBench therefore includes a blinded, offline,
right-to-left review application. A release bundle contains no model
prediction. Three qualified native reviewers independently assign the full
label set, a three-level confidence judgment, and optional flags for ambiguity,
possible bad gold, personally identifying information, and distressing
content. The ``not'' label is enforced as mutually exclusive.

Aggregation validates task IDs and label sets, requires complete submissions,
computes unanimous and majority consensus, and surfaces disagreement for
adjudication. Reviewer IDs are assigned pseudonyms; progress remains in local
browser storage until an explicit JSON export. This protocol is prepared for a
300-item AraTox audit. Human results are intentionally not fabricated in this
technical report.

\section{What would justify annotator replacement?}
The v0.1 evidence is deliberately below the bar for replacing human
annotators. A credible replacement or auto-acceptance claim would require at
least four additions.
\begin{enumerate}
  \item \textbf{Human reference study.} Run three independent qualified native
  reviewers on the same blind bundle, report Krippendorff's $\alpha$ or an
  appropriate multi-label agreement statistic, estimate gold-label error, and
  adjudicate persistent disagreements before treating the dataset as a
  replacement target.
  \item \textbf{Frontier and specialist snapshots.} Evaluate dated snapshots
  of the systems likely to be deployed: frontier APIs, larger Arabic-specialist
  models, and a fine-tuned Arabic encoder such as MARBERT-style baselines
  where licensing permits. Local sub-10B quantized runs are useful smoke tests,
  not production evidence.
  \item \textbf{Prompt and calibration matrix.} Cross zero-shot and few-shot
  prompts, English and Arabic instructions, self-consistency, and confidence
  elicitation. Selective automation should then be evaluated on a held-out
  calibration/evaluation split and compared with alt-test or conformal
  risk-control procedures.
  \item \textbf{Power and economics.} Increase paired sample sizes before
  turning small differences into claims; report latency, cost per 10,000
  attempted labels, cost per correctly accepted label, and the share escalated
  to humans.
\end{enumerate}
Under that protocol, the central output would not be a single accuracy table.
It would be a decision curve: expected error rate and cost as automation
coverage increases, with human review used wherever the quality floor is not
statistically supported.

\section{Limitations and ethics}
This is a launch audit, not a top-tier benchmark paper or ranking of all
frontier systems. The principal generative results are local quantized models;
closed frontier systems and larger Arabic-specialist models remain untested.
The prompts are zero-shot and English-authored; no few-shot, Arabic-prompt,
chain-of-thought, self-consistency, or prompt-sensitivity axis is reported.
The headline LLM samples are modest (100, 300, and 500 items), while the
supervised diagnostics train in-domain and therefore should be interpreted as
learnability controls rather than competitive model baselines.

The human adjudication protocol is implemented but not yet run, so v0.1 cannot
estimate gold-label ambiguity, inter-annotator agreement, or human replacement
validity. Future releases should add three-rater adjudication with
Krippendorff's $\alpha$, dated frontier snapshots, native Arabic and few-shot
prompts, NADI/OSACT-compatible tasks, held-out calibration splits, and an
alt-test or conformal risk-control analysis before making annotator-replacement
claims.

The dialect evaluation covers broad regions, not countries, cities, register,
code-switching, Arabizi, or speech. Dialect is socially variable and cannot be
reduced to geography. DART itself reported lower human agreement for Iraqi and
Maghrebi labels \citep{dart2018}; model errors in those groups must not be
interpreted as deficiencies in the varieties or their speakers.

Toxicity data can harm reviewers. The adjudication interface includes a
distressing-content flag, supports offline work, and should be deployed with
informed consent, opt-out, breaks, and access to support. Public release must
follow each source's terms. LabelBench redistributes no IADD/DART text and
publishes only aggregate metrics and cryptographic manifest identifiers.

\section{Conclusion}
Arabic labeling automation cannot be certified by a binary accuracy number.
Across toxicity and dialect tasks, the work becomes sharply harder when the
label space reflects production reality. The correct question is not ``Can a
model label Arabic?'' but ``Which taxonomy, which variety, at what proven
quality floor, and with how much human review?'' LabelBench makes that question
measurable.

\bibliographystyle{plainnat}
\bibliography{references}
\end{document}
